\section{实际执行中的差异}\label{sec:execution}

前几章的比较都在信号层与统计层。但业界与文献反复给出同一个教训：\emph{统计上显著不等于能赚钱}。Martin \& Sch\"oneborn(2011) 把这个教训写进了标题——“均值回归能赚，但要付代价”\cite{martin2011}。本章把第~\ref{sec:duality}~节的对偶翻译成执行层面的具体差异，按从信号到资金的顺序展开。

\subsection{第一步就分岔：信号平滑与换手控制}

两个策略在离开信号层的第一步就走了不同的路。

\textbf{趋势一侧的核心工程问题是“少交易”}。动量信号在日频上噪声极大，若直接按日频符号调仓，换手会高到成本吞掉全部收益。文献给出的主流解法有三类：一是用线性趋势拟合的斜率替代离散符号（Baltas \& Kosowski 2013 报告线性趋势拟合信号样本外最优且换手最低\cite{baltaskosowski2013imp}）；二是用区间估计量替代收盘价波动率以降低噪声（Yang--Zhang 估计量在效率与偏差间取得平衡，同文献）；三是对头寸做指数平滑或部分调整（Gârleanu--Pedersen 的“目标前置、部分调整”给出了含成本的动态最优解\cite{garleanu2013}）。Kim 等(2016) 的系统检验进一步指出，时间序列动量的收益中有相当部分来自波动率缩放本身而非时序符号\cite{kim2016}——这意味着平滑与缩放的技术选择会显著改变策略的经济含义。

\textbf{反转一侧的核心工程问题是“带宽”}。反转策略不能频繁穿越阈值，否则会在价差尚未收敛时反复止损。理论给出的答案是：最优交易带（no-trade band）的宽度随交易成本以三次方根增长\cite{martin2011}。这一结论的实操含义是：\emph{成本越高的市场，反转策略应当越“迟钝”}，宁愿漏掉一部分信号，也不要在噪声中频繁进出。Boguslavsky 与 Boguslavskaya(2004) 在幂效用下给出的解析阈值\cite{boguslavsky2004}、Zhang 与 Zhang(2008) 的“低买高卖”最优规则\cite{zhangzhang2008}、以及 Leung 与 Li(2015) 引入止损退出的双停时解\cite{leungli2015}，构成了这一族的技术基础。

两者的对称性值得注意：趋势在“平滑信号”，反转在“拓宽带宽”。前者压的是信号的方差，后者压的是交易的频率。目标函数都是换手成本，手段却因信号的可加性不同而不同——趋势信号可以连续加权，反转信号必须通过阈值离散化。

\subsection{组合构建与中性化}

在两个策略上，组合层的任务不同。

\textbf{趋势策略需要“去相关”与“风险平价”}。跨资产的趋势信号在危机中高度同向，因此组合必须按资产波动率与相关性分配风险，否则单一资产会主导组合。行业/板块层面的中性化在股票动量中是必需的——Moskowitz \& Grinblatt(1999) 指出行业动量解释了大部分个股动量\cite{moskowitzgrinblatt1999}，若不中性化，组合实际上在押注行业轮动而非个股趋势。

\textbf{反转策略需要“彻底的非系统性”与“分散”}。配对交易的价差交易天然对冲了共同因子暴露，这是它的优势；但代价是\emph{特质风险无法被市场因子抵消}，因此必须靠配对数量分散。Avellaneda \& Lee(2010) 的 PCA 残差框架正是为这一目的设计的：先剥离前若干主成分（共同因子），再在残差上做均值回复\cite{avellaneda2010}。若不做因子剥离，反转组合会混入大量方向性风险，其“收敛”叙事就不成立。

一个容易被忽略的工程细节：反转策略的\emph{对冲比率}是估计量而非常数。Elliott 等(2005) 用 Kalman 滤波在线估计配对价差的 OU 状态\cite{elliott2005}，其工程含义是：对冲比率随时间变化，若沿用静态比率（如形成期的回归系数），价差序列会因参数漂移而失真。趋势策略没有对应的困扰，因为它的仓位由自身波动率决定，不依赖跨资产的稳定关系。

\subsection{波动率目标化的不对称作用}

这是本文认为执行层最重要的一条差异。

波动率目标化的做法是：按“目标波动率 / 已实现波动率”缩放仓位。它在趋势策略上被反复证明有效——Barroso \& Santa-Clara(2015) 把动量夏普从 0.53 提到 0.97、超额峰度从 18.2 降到 2.7\cite{barroso2015}；Harvey 等(2018) 在 60 多个资产上系统验证其改善风险资产夏普\cite{harvey2018}；Moreira \& Muir(2017) 表明波动率管理能提升多数因子的表现\cite{moreiramuir2017}。

但把这套方法搬到反转策略上会产生\emph{方向性错误}。理由是信号与波动率的关系在两策略中符号相反：趋势策略在波动率上升时信号质量下降（趋势被噪声打断），降敞口是保护；反转策略的利润经常\emph{恰好}来自高波动时段——偏离更大，回归的绝对空间也更大。因此在波动率上升时降低反转敞口，等于在信号最强时把仓位压到最小。

正确做法是把波动率信息用在\emph{风险预算}而非\emph{信号缩放}上：对反转策略，波动率上升应当触发的是止损线的收紧与资金冗余的检查，而不是机械的仓位缩放。Leung \& Li(2015) 把止损退出纳入最优框架，正是这一思路的理论化\cite{leungli2015}。

\subsection{成本与执行机制}

成本模型的选择会系统性影响两策略的比较结论。

\textbf{趋势策略的成本压力来自再平衡}。月频或周频调仓下，单次调仓的换手中等，但一年累积的换手仍可观；成本主要在“定期滚动”中流失，可以用信号平滑直接压低。

\textbf{反转策略的成本压力来自进出频率}。它的成本具有明显的非线性：价差偏离越大时，往往也是流动性越差的时刻（流动性螺旋），因此实际冲击成本高于按平均换手估算的值。Brunnermeier \& Pedersen(2009) 的模型给出了这一非线性\cite{brunnermeier2009}。工程上应当采用\emph{参与率约束的冲击成本模型}（冲击与参与率的平方根成正比），并对超过参与率上限的部分显式标记、不进入回测收益。

执行算法层面，Almgren \& Chriss(2000) 的框架给出冲击成本与择时风险的权衡\cite{almgrenchriss2000}；Perold(1988) 的“执行缺口”提供了一致的度量口径\cite{perold1988}。对反转策略，执行缺口的斜率更陡（因为它常在流动性差时交易）；对趋势策略，缺口更平坦但累积量更大。

\subsection{资金、杠杆与失效防线}

第~\ref{sec:duality}~节的对偶六指出两策略的失效模式不同，执行层的应对也必须不同。

反转策略的防线是\emph{资金冗余}。Shleifer \& Vishny(1997) 的套利限制理论说明，由于业绩与资金流正反馈，套利者在错价加深时反而被迫减仓\cite{shleifervishny1997}；2007 年 8 月的事件表明这一机制在大规模、多机构同时发生时会被放大\cite{khandanilo2011}。因此反转策略必须预留足以承受“偏离继续扩大”的保证金缓冲，并设定明确的偏离容忍上限——超过上限即减仓，而不是“越跌越买”。

趋势策略的防线是\emph{状态过滤}。动量崩溃发生在恐慌状态（市场下跌后、波动率高企并伴随反弹），且部分可预测\cite{danielmoskowitz2016}。因此趋势策略需要在极端状态下降低甚至反转敞口，而不是无条件追随短期走势。Chabot 等(2014) 的历史研究也指出动量与“追涨”行为在可预测的崩溃前会呈现特征性模式\cite{chabot2014}。

\subsection{Regime 依赖与切换}

两策略的适用状态不同，这一点在文献中有明确的结构。

趋势策略在\emph{持续单向运动}的环境中有效，这类环境通常伴随宏观变量的持续漂移；反转策略在\emph{区间震荡}环境中有效。Hamilton(1989) 的马尔可夫切换模型提供了建模基础\cite{hamilton1989}，Ang \& Bekaert(2002) 把它用于国际资产配置并显示相关性本身随状态变化\cite{angbekaert2002}，Ang \& Timmermann(2012) 的综述系统梳理了 regime 方法在金融中的应用\cite{angtimmermann2012}。

工程上，regime 的用途有两种。一是\emph{硬切换}：只在某状态下启用某策略；二是\emph{软加权}：按状态概率调整两策略的权重。Wood 等(2021) 的“慢动量加快反转”属于更彻底的方案——用变点检测识别状态转换点，让高频分量做反转、低频分量做趋势\cite{wood2021}。这一方案的动机与本文统一框架一致：既然两者是同一信号的不同频段，就不必在它们之间做二选一。

\subsection{组合融合：为什么简单叠加常常不分散}

一个常见的组合做法是把趋势与反转等权叠加，期望获得分散化。这在多数时候不起作用，原因在第~\ref{sec:mechanism}~节已经点出：两策略在某些状态下相关性趋近于一（共同的资金冲击），而在另一些状态下相关性为负（信号符号相反）。因此简单叠加的实际效果是\emph{在信号相反时互相抵消，在资金冲击时同时亏损}——两头不讨好。

更合理的做法有三条。一是\emph{按尺度分层}：短尺度配反转、中尺度配趋势，让不同频段各自发挥作用（对应 Wood 等的频段分解）。二是\emph{按状态加权}：用 regime 概率动态调整权重，而非固定比例。三是\emph{按风险预算而非资金分配}：先评估两策略在目标状态下的波动与尾部，再分配风险预算，避免某一策略在极端状态下主导组合。

\subsection{执行差异清单}

\begin{table}[htbp]
\centering\footnotesize
\caption{执行层面的差异清单}\label{tab:exec}
\begin{tabular}{@{}>{\raggedright\arraybackslash}p{3.0cm}>{\raggedright\arraybackslash}p{5.9cm}>{\raggedright\arraybackslash}p{5.9cm}@{}}
\toprule
\textbf{执行环节} & \textbf{趋势跟踪} & \textbf{均值回归} \\
\midrule
信号处理 & 平滑信号、压低换手；可用斜率替代符号 & 拓宽交易带（宽度 $\propto$ 成本立方根） \\
对冲设计 & 不依赖跨资产关系 & 依赖对冲比率，需在线估计（Kalman） \\
因子暴露 & 需行业/板块中性化，去同向相关 & 需彻底剥离共同因子，纯特质化 \\
风险缩放 & 波动率目标化有效 & 波动率目标化可能有害，改用止损+资金冗余 \\
成本结构 & 累积式，可由平滑压低 & 非线性式，流动性差时冲击更大 \\
资金防线 & 状态过滤，防动量崩溃 & 资金冗余 + 偏离上限，防强制平仓 \\
适用状态 & 持续单向运动、宏观漂移 & 区间震荡、偏离扩大 \\
组合融合 & 需按尺度或状态加权 & 简单叠加会互相抵消 \\
\bottomrule
\end{tabular}
\end{table}

这张清单是本文对实践最直接的一份交付。它给出的核心信息是一句反直觉的话：\textbf{两策略在信号上越接近，执行上的区别就越关键}——因为信号带来的边际 alpha 有限，而执行偏差（成本、对冲、资金）足以决定正负。
